\subsection{Архитектура предлагаемого решения}

1. сначала искал готовые хранилища obsidian. не важно, есть в них ссылки или нет. 
достаточно было найти базу из заметок, чтобы потом разметить.
по результатам нашел несколько готовых проработанных хранилищ:
https://github.com/noodleslove/notes/
https://github.com/hundredrabbits/100r.co
https://github.com/lukesmithxyz/based.cooking 
https://github.com/keyvanakbary/learning-notes


изучив их подробнее, понял, что они не подходят





На начальном этапе проектирования системы стояла задача подбора подходящего набора данных для последующей разметки и использования в качестве валидационного множества («золотого стандарта») для оценки качества работы GraphRAG. Изначальной гипотезой был поиск готовых пользовательских хранилищ (vaults) формата Obsidian. Наличие или отсутствие явных связей в них на данном этапе не являлось критичным фактором, так как основной целью был поиск репрезентативной текстовой базы для последующего ручного извлечения и семантической классификации связей.

В ходе предварительного анализа были рассмотрены несколько проработанных open-source хранилищ, размещенных на платформе GitHub:

    noodleslove/notes — обширная база личных заметок.

    hundredrabbits/100r.co — структурированная документация и база знаний цифрового коллектива.

    lukesmithxyz/based.cooking — репозиторий кулинарных рецептов.

    keyvanakbary/learning-notes — база знаний, сфокусированная на разработке программного обеспечения и инженерии.

Несмотря на открытость и доступность данных форматов, детальный анализ выявил ряд существенных недостатков, делающих эти хранилища непригодными для решения задачи валидации сложных графовых связей. Выделены следующие ключевые проблемы:

    Семантическая разреженность связей: Значительная часть заметок (например, в репозитории кулинарных рецептов) имеет тривиальную структуру. В них отсутствуют нетривиальные, многоуровневые логические взаимосвязи (например, причинно-следственные, иерархические или концептуальные), что делает бессмысленным применение сложных алгоритмов классификации отношений.

    Проблема масштабирования и фрагментации: Рассмотренные базы обладают слишком большим объемом для проведения качественной ручной разметки «золотого стандарта» в рамках данного исследования. Изоляция отдельного фрагмента (сабграфа) неизбежно приведет к потере глобального контекста и повисшим ссылкам (dangling links), что исказит метрики RAG-системы.

    Низкая информационная плотность и качество контента: Содержание многих личных хранилищ специфично: зачастую это обрывочные предложения, неструктурированные куски программного кода, цитаты без контекста или прямая «копипаста» из сторонних статей. Подобный контент плохо подходит для обработки языковыми моделями (LLM), требующими связного повествовательного текста для качественного извлечения сущностей и связей.

    Лицензионные и этические ограничения: Большинство персональных баз знаний публикуются без четкого указания лицензии, допускающей модификацию и создание производных наборов данных для машинного обучения, что создает потенциальные риски при публикации результатов исследования.

На основании вышеизложенного было принято решение отказаться от использования пользовательских хранилищ Obsidian в пользу корпуса текстов свободной энциклопедии Википедия. Данный выбор обоснован следующими преимуществами:

    Высокое качество и связность текста: Статьи написаны в строгом энциклопедическом стиле, обладают высокой информационной плотностью и четкой структурой, что оптимально для NLP-задач.

    Наличие размеченной структуры: Википедия уже содержит богатую сеть гиперссылок и структурированные данные (инфобоксы), которые могут служить отличным базовым ориентиром (baseline) или источником слабого контроля (weak supervision) при автоматической валидации связей.

    Удобство автоматической обработки: Наличие стандартизированных XML-дампов и открытого API позволяет легко парсить контент, гибко фильтровать статьи по предметным областям и извлекать только релевантные текстовые блоки, формируя репрезентативный набор данных контролируемого объема.

Пара советов по тексту:

    Если научрук любит конкретику, в финальной версии диплома в блоке про Википедию можно будет добавить 1-2 предложения о том, как именно ты ее парсишь (например, используешь библиотеку wikipedia-api или качаешь конкретный дамп).

    Про лицензии — это очень сильный аргумент для академической работы, он показывает зрелый подход к сбору данных. Обязательно оставь его
